Durante el desarrollo del presente proyecto se han identificado diversas oportunidades de mejora y líneas de evolución del sistema que no han sido abordadas en esta primera versión, pero que resultan coherentes con los objetivos y el alcance definidos en la introducción.

En particular, el sistema ha sido concebido con un enfoque modular y extensible, lo que permite la incorporación progresiva de nuevas funcionalidades y la mejora de las ya existentes sin necesidad de rediseñar la arquitectura base. Las líneas de trabajo futuro que se presentan a continuación se organizan en función de los módulos funcionales definidos en el sistema, reflejando así su evolución natural.

\subsection*{M1. Gestión de acceso, alta y autorización de usuarios}
\begin{itemize}
    \item Implementación de mecanismos de recuperación de contraseña y verificación de cuentas de usuario.
    \item Incorporación de medidas adicionales de seguridad frente a accesos indebidos, como la limitación de intentos de autenticación.
    \item Gestión avanzada de sesiones, incluyendo el control de sesiones activas y su posible revocación.
    \item Desarrollo de herramientas específicas para la visualización y análisis de los registros de acceso y acciones administrativas.
\end{itemize}